\section{Two Vocabularies, and What Survives Composition}
\label{sec:ch15-two-vocabularies}
\index{authorization!the two families of attribute|(}%

Here is \code{Customer}, as Accounts declares it at this tag. Look at the last
field and count the attributes:

\begin{minted}[fontsize=\footnotesize,breaklines]{csharp}
[Key("id")]
public sealed record Customer(
    Guid Id,
    string DisplayName,
    [property: Authenticated]
    [property: Authorize]
    string Email)
\end{minted}

Two attributes that read like synonyms. They come from different packages, they
are enforced by different processes, and neither one does the other's job.

\index{Authenticated@\code{[Authenticated]}!emits and does not enforce}%
The first comes from the federation package,
\code{HotChocolate.ApolloFederation}. It prints \code{@authenticated} into this
service's published schema, the composer copies it into the router's field
configuration, and the router refuses the field to a caller with no token.
Inside this process it does nothing whatsoever.

\index{Authorize@\code{[Authorize]}!enforces and does not compose}%
The second comes from \code{HotChocolate.Authorization}, which Mosaic gets
through the ASP.NET Core integration package that wires it to a policy engine.
It puts middleware in front of the resolver in this service. It knows nothing
about federation, and the composer, as we are about to see, throws it away.

\subsection*{Why the first one cannot enforce anything}

\index{HotChocolate.ApolloFederation@\code{HotChocolate.ApolloFederation}!what it references}%
I expected to establish this with a grep and found something better. The
federation package's own project file declares exactly one project reference:

\begin{minted}[fontsize=\footnotesize,breaklines]{text}
<ProjectReference Include="..\..\..\Core\src\Core\HotChocolate.Core.csproj" />
\end{minted}

\code{HotChocolate.Authorization} is a different package and is not among them,
so there is no path from \code{[Authenticated]} to any authorization handler.
This is not a gap somebody forgot to close. The attributes are schema
annotations, they are documented as instructions to a supergraph router, and
the package that owns them could not enforce them without taking a dependency it
does not have. Read at tag \code{16.6.0}, commit \code{8fea46e}; the three
attributes and the directive types they configure are in
\code{Types/Directives/}, and it is the directive types that carry the
\code{[Package(...)]} version gate, at federation v2.5 for \code{@authenticated}
and \code{@requiresScopes} and v2.6 for \code{@policy}.

The consequence is the thing to carry. A subgraph that declares
\code{@requiresScopes} and nothing else is protected only for as long as every
caller arrives through the router. For a graph whose services sit on a private
network that may be a fair bet. For Mosaic, whose seven services publish host
ports so that a reader can \code{curl} them, it is not a bet at all, and that is
why \code{Customer.email} carries two attributes rather than one.

\subsection*{What the composer keeps}

\index{composition!the authorization field configuration}%
Compose the eight subgraphs and the routing table gains something
chapters~\ref{ch:composition} and~\ref{ch:enter-the-router} never saw:

\begin{minted}[fontsize=\footnotesize,breaklines]{json}
{
  "typeName": "Query",
  "fieldName": "customerById",
  "authorizationConfiguration": {
    "requiresAuthentication": true,
    "requiredOrScopes": [{"requiredAndScopes": ["customers:read"]}]
  }
}
\end{minted}

Four fields carry one at this tag. Two details in that fragment are worth more
than a glance.

\index{requiresScopes@\code{@requiresScopes}!implies authentication}%
The first is that \code{customerById} never said \code{@authenticated}
anywhere. It carries \code{[RequiresScopes]} alone, and the composer wrote
\code{requiresAuthentication: true} as well, because a caller who must hold a
scope must first be somebody. One directive, two flags, and the consequence is
that a schema review counting \code{@authenticated} annotations will undercount
what the graph enforces.

\index{requiresScopes@\code{@requiresScopes}!the shape of its argument}%
The second is \code{requiredOrScopes} wrapping \code{requiredAndScopes}. The
directive's argument is a list of lists: the outer list is an OR and the inner
list an AND, so one scope is written \code{[["customers:read"]]} and
\code{[["a","b"],["c"]]} means \code{a} and \code{b}, or else \code{c}. That
reads oddly and it is the specification's shape rather than HotChocolate's
invention. Applying \code{[RequiresScopes]} twice to one field appends a second
OR branch rather than replacing the first.

\subsection*{Three names for one scalar}

\index{openfed\_\_Scope@\code{openfed\_\_Scope}}%
Now introspect the router. The client-facing schema declares both directives,
which is already interesting: chapter~\ref{ch:composition} found that a Cosmo
execution config carries no \code{join\_\_} directives at all, and these two
survive. What it declares is not quite what the subgraph published:

\begin{minted}[fontsize=\scriptsize,breaklines]{text}
subgraph:  directive @requiresScopes(scopes: [[String!]!]!) on ...
composed:  directive @requiresScopes(scopes: [[openfed__Scope!]!]!) on ...
           scalar openfed__Scope
\end{minted}

Three projects, three names for it. Apollo's reference says
\code{federation\_\_Scope}~\autocite{apollo2026authdirectives}. Cosmo says
\code{openfed\_\_Scope}. HotChocolate declares a \code{Scope} scalar of its
own, then overrides the directive's argument to plain \code{String} with a
\code{[GraphQLType]} attribute, so what a HotChocolate subgraph publishes is
neither of those. The composer quietly substitutes its own and declares the
scalar, so that the document still makes sense.

Nothing breaks, because all three serialise as strings. What it costs is that
\code{scalar openfed\_\_Scope} is now in Mosaic's published schema, visible to
anyone who introspects the router, chosen by nobody. It has company:
\code{scalar FieldSet} has been leaking the same way since
chapter~\ref{ch:composition}, and chapter~\ref{ch:security-hardening} owns
deciding whether a production graph does anything about either.

\subsection*{What the composer throws away, and what it keeps by accident}

\index{authorize@\code{@authorize}!printed, then dropped}%
So much for the federation pair. Here is what happens to the other one.

\code{@authorize} is printed in full into the SDL a composer reads. That
surprised me, because its directive type calls \code{.Internal()}, and I had
assumed that meant hidden. It does mean hidden, from
\code{\{ \_\_schema \{ directives \{ name \} \} \}}, because \code{IsPublic} is
consulted when building the introspection listing. The SDL printer behind
\code{\_service \{ sdl \}} does not consult it, and \code{\_service \{ sdl \}}
is the field a federation composer actually reads. So the committed
\code{schema/accounts.graphql} carries this:

\begin{graphqlsdl}[fontsize=\footnotesize]
type Query {
  customerById(id: ID!): Customer
    @authorize(policy: "CustomersRead")
    @requiresScopes(scopes: [["customers:read"]])
  ...
}

type Customer implements Node @key(fields: "id") {
  id: ID!
  displayName: String!
  email: String! @authenticated @authorize
}
\end{graphqlsdl}

together with a full \code{directive @authorize(...)} definition, because a
subgraph schema is a document that has to stand on its own.

Compose it and \code{@authorize} is gone, without being rejected or warned
about: there are zero occurrences of the string in the client-facing schema, and
no authorization configuration anywhere that came from it. That is correct
behaviour. It is not a federation directive, the composer has never heard of it,
and dropping an unknown directive is what a composer should do with one.

What it does not drop is this:

\begin{minted}[fontsize=\footnotesize,breaklines]{text}
"""Defines when a policy shall be executed."""
enum ApplyPolicy {
  """Before the resolver was executed."""
  BEFORE_RESOLVER
  """After the resolver was executed."""
  AFTER_RESOLVER
  """The policy is applied in the validation step before the execution."""
  VALIDATION
}
\end{minted}

\index{ApplyPolicy@\code{ApplyPolicy}!in the published graph}%
That enum is in Mosaic's published client schema. It is there because it is the
type of \code{@authorize}'s \code{apply} argument, and the directive that gave
it a reason to exist was removed while it was not. Anybody who introspects this
graph, or generates a client from it, is told about the resolver execution
phases of one subgraph's server library.

I am confident about where it came from rather than guessing, because the case
is committed both ways: strip every \code{@authorize} usage and both definitions
out of the three subgraphs that carry them, recompose, and the authorization
table is identical while \code{ApplyPolicy} disappears. Two things fall out of
that. The enum is the directive's fault, and a subgraph that stops enforcing
anything locally composes to exactly the same routing table as one that does,
which is the whole problem with having the rule written twice.

\subsection*{The directive that composes and does nothing}

\index{policy@\code{@policy}!accepted and discarded}%
There is a third federation directive, and this is the part of the chapter I
would most like to be wrong about.

\code{@policy} is in Apollo's federation directive
reference~\autocite{apollo2026authdirectives}, beside the two that work.
HotChocolate ships a \code{[Policy]} attribute that emits it, gated at
federation v2.6, which is the version this stack composes at by default. So a
subgraph author can reach for it with no ceremony at all.

So annotate a field with it. Inventory's \code{availableQuantity} gets this,
with the \code{@link} import and the directive definition beside it:

\begin{graphqlsdl}[fontsize=\footnotesize]
availableQuantity: Int! @policy(policies: [["stock:read"]])
\end{graphqlsdl}

Compose, and wgc 0.129.7 exits zero and prints its ordinary success line. The
composed config has no authorization configuration for that field. Not an empty
one: none. The client schema has no \code{@policy} in it. There is no warning
and no error.

The reason is in the composer's source. Version 0.63.2 of
\code{@wundergraph/composition} defines \code{@authenticated} and
\code{@requiresScopes} in its directive table and has no \code{@policy} entry of
any kind, under that name or Cosmo's own. The router agrees: its authorizer
enforces the two and mentions no policy.

So the sequence is: a directive the specification defines, which the subgraph
library ships an attribute for, annotated onto a field by somebody who has read
the specification, accepted by the composer, and enforced by nothing. There is
no step in that chain where anyone is told.

Chapter~\ref{ch:real-time-in-a-federated-world} ended on a directive the router
advertises and does not honour, and said that a directive in a published schema
is a promise. This is the same failure with the polarity reversed and the stakes
raised. \code{@defer} promises a client something the graph will not do, and
costs a loading state. \code{@policy} promises a \emph{developer} something the
graph will not do, and costs whatever was behind the field. Both are committed
as cases now, so a release that fixes either fails the gate rather than making a
chapter quietly wrong.

\medskip
That settles what the router knows. The next question is what reaches the
service behind it, because two of Mosaic's three layers live there and neither
can do anything without a caller.

\index{authorization!the two families of attribute|)}%
